\chapter{Хранение транспортера-тягача}
\section{Общие положения}

Порядок постановки транспортера-тягача на~хранение, последовательность и~организация работ, выполняемых при~подготовке и~содержании транспортера-тягача на~хранение, а~также порядок снятия его с~хранения определяется Наставлением по~автотракторной службе Вооруженных Сил РФ и~Инструкцией по~хранению автотракторной техники и~имущества в~воинских частях, на~базах и~складах. 

Постановке на~хранение подлежат все транспортеры-тягачи, эксплуатация которых не~планируется на~срок более трех месяцев, а~в~особых климатических условиях~--- более одного месяца. 

В~зависимости от~срока хранения транспортера-тягача она может быть кратковременной (на~срок до~одного года) и~длительный (на~срок более одного года). 

Транспортер-тягач, подлежащий кратковременному хранению, подвергается очередному номерному техническому обслуживанию (ТО-1 или~ТО-2), при~проведении которого выполняется ряд специальных дополнительных операций по~подготовке его деталей, механизмов и~агрегатов к~консервации. 
\looseness=-1

Транспортер-тягач, подлежащий длительной хранению, подвергается техническому обслуживанию №2, при~проведении которого необходимо выполнить все работы, предусмотренные через одно ТО-2 и~при~сезонном обслуживании. Одновременно с~техническим обслуживанием выполняются дополнительные работы по~подготовке транспортера-тягача к~хранению. 
\looseness=-1

Контроль за~состоянием транспортеров-тягачей, находящихся на~хранении, и~их техническое обслуживание проводятся в~следующие сроки: 
\begin{itemize}
    \item один раз в~месяц (в~парковые дни);
    \item один раз в~шесть месяцев (при~подготовке транспортеров-тягачей к~летнему и~зимнему периодам эксплуатации); 
    \item один раз в~год (в~сухую, без~осадков, погоду).
\end{itemize}

\begin{table}[h]
    \center
    \begin{tabular}{|m{2cm}|>{\tc}m{2.5cm}|>{\tc}m{2.5cm}|>{\tc}m{2.5cm}|}
        \hline
        \multicolumn{2}{|m{4.5cm}|}{\multirow{2}{4.5cm}{\tc Вид трудозатрат}} & \multicolumn{2}{m{5cm}|}{\tc Трудозатраты на~подготовку к~консервации, \textit{чел.-час} } \\
        \cline{3-4}
        \multicolumn{2}{|c|}{} & кратко\-временную  & длительную \\
        \hline
        \multicolumn{2}{|m{4.5cm}|}{Без~учета ТО-1 и~ТО-2}  & 13,5 & 18,2 \\
        \hline
        \multirow{2}{2cm}{\centering С учетом ТО-1 и~ТО-2} & без~регулировочных работ & \range{20,5}{26,5} & \range{40,2}{54,5} \\
        \cline{2-4}
        & с~регулировочными работами & \range{23,1}{28,5} & \range{44,9}{57,2} \\
        \hline

    \end{tabular}
    \caption{Трудозатраты на~подготовку к~консервации}
    \label{worktime_conserv}
\end{table}

Кроме этого, ежегодно опробуется 50\% транспортеров-тягачей, находящихся на~длительном хранении, из~них: 30\% проверяется с~пуском двигателя и~прокручиванием механизмов силовой передачи на~месте; 20\%~--- контрольным пробегом на~15~км, с~последующей полной переконсервацией. 

Таким образом за~5 лет хранения каждый транспортер-тягач должен подвергаться контрольному пробегу. 

В~первый год постановки транспортера-тягача на~длительное хранение, опробование его с~запуском двигателя не~производится.

\begin{table}[t]
    \begin{center}
        \begin{tabular}{|m{2.5cm}|m{3cm}|>{\tc}m{2cm}|>{\tc}m{2cm}|}
            \hline
            \multicolumn{2}{|m{5.5cm}}{\multirow{2}{5.2cm}{\tc Периодичность технического обслуживания}} & \multicolumn{2}{|>{\tc}m{4cm}|}{Трудозатраты при~консервации, \textit{чел.-час}} \\
            \cline{3-4}
            \multicolumn{2}{|m{5.5cm}|}{} & кратко\-временной & длительной \\
            \hline
            \multicolumn{2}{|m{5.5cm}|}{Один раз в~месяц}  & \range{0,9}{1} & \range{0,9}{1}\\
            \hline
            \multicolumn{2}{|m{5.5cm}|}{Один раз в~шесть месяцев}  & \range{3,8}{4} & \range{3,8}{4} \\
            \hline
            \multicolumn{2}{|m{5.5cm}|}{Один раз в~год}  &~--- & \range{4,6}{5} \\
            \hline
            \multirow{2}{2.5cm}{Опробование транспортера тягача:} & с~пуском двигателя и~прокручиванием силовой передачи на~месте  &~--- & \range{19}{19,5}\\
            \cline{2-4}
            & контрольным пробегом&~--- & \range{20,8}{21,3} \\
            \hline

        \end{tabular}
        \caption{Трудозатраты на~техническое обслуживание и~опробование транспортера-тягача, находящегося в~консервации}
        \label{worktime_obsluz}
    \end{center}
    \note{Трудозатраты на~опробование транспортера-тягача даны с~учетом работ на~снятие с~хранения, опробования и~постановку на~хранение.}
\end{table}

Работы по~обслуживанию и~опробованию транспортеров-тягачей, находящихся на~длительном хранении, должны проводиться в~соответствии с~календарным планом-графиком, разрабатываемым в~каждой части на~5 лет. 

Техническое обслуживание, проведенное один раз в~шесть месяцев и~один раз в~год, а~также опробование транспортера-тягача записываются в~специальном вкладыше его паспорта. 

Быстрое приведение транспортера-тягача, находящегося на~хранении, в~боевую готовность обеспечивается: 
\begin{itemize}
    \item точностью выполнения всех операций по~его подготовке к~консервации и~по~снятию с~консервации согласно рекомендуемой технологии; 
    \item знанием и~правильным выбором наиболее эффективных способов пуска и~прогрева двигателя; 
    \item выбором рационального способа хранения аккумуляторных батарей и~четкой организацией работ, обеспечивающих приведение батарей в~рабочее состояние с~минимальными затратами времени; 
    \item выполнение строго ограниченного перечня работ первой очереди, позволяющих пустить двигатель и~вывести транспортер-тягач из~парка с~таким расчетом, чтобы работы второй очереди выполнить при~первой возможности в~пути (в~районе сосредоточения, на~привалах и~остановках). 
\end{itemize}

\begin{table}[h]
    \begin{center}
        \begin{tabular}{|l|c|>{\tc}m{2.1cm}|>{\tc}m{2.1cm}|}
            \hline
            \multicolumn{2}{|c|}{\multirow{2}{*}{Вид работ}} & \multicolumn{2}{>{\tc}m{4.4cm}|}{Продолжительность работ при~расконсервации, \textit{мин}} \\
            \cline{3-4}
            \multicolumn{2}{|c|}{} & кратко\-временной & длительной \\
            \hline
            \multirow{2}{*}{Работы первой очереди} & летом  & 15 & 20 \\
            \cline{2-4}
                                                    & зимой  & 30 & 35 \\
            \hline
            \multicolumn{2}{|l|}{Работы второй очереди}     & 20 & 20 \\
            \hline
            \multirow{2}{*}{Общее время}           & летом  & 35 & 40 \\
            \cline{2-4}
                                                    & зимой  & 50 & 55 \\
            \hline

        \end{tabular}
        \caption{Средняя продолжительность работ при~снятии транспортера-тягача с~консервации}
        \label{average_wortime}
    \end{center}
    \note{Средняя продолжительность работ дана при~условии хранения аккумуляторных батарей в~рабочем состоянии: летом~--- на~транспортерах-тягачах, зимой~---\\в~аккумуляторной (без~учета времени на~доставку их к~машине).}
\end{table}



\section{Подготовка к~кратковременному хранению}
\subsubsection{Содержание работ и~технические условия}
\begin{enumerate}
    \item \begin{enumerate}
        \item Тщательно очистить от~грязи и~вымыть транспортер-тягач, слить воду из~корпуса, протереть насухо наружные поверхности агрегатов и~узлов, из~труднодоступных мест воду удалить сжатым воздухом, 
        
        \item При~проведении уборочно-моечных работ следить, чтобы не~попадала вода в~впускную и~выпускную системы двигателя, котел и~редуктор подогревателя и~на~приборы электрооборудования. 
            
        \item Выполнить очередное номерное техническое обслуживание (ТО-1 или~ТО-2). 
    \end{enumerate}

    \item \begin{enumerate}
        \item Промыть, протереть и~осмотреть дюритовые шланги. Шланги, имеющие расслоения с~торцов, вздутия и~широкие (до~1~мм) трещины на~поверхности, заменить. 
        \item Очистить и~окрасить стяжные хомуты шлангов. Резьбовую часть стяжных болтов смазать консервационной смазкой СХК или~ПВК. 
    \end{enumerate}
    

    \item Заправить систему охлаждения: летом~--- водой с~трехкомпонентной присадкой (по. 0,05\% весовых нитрита натрия, тринатрия фосфата и~двухромокислого калия); зимой~--- низкозамерзающей охлаждающей жидкостью (антифризом) с~этой~же присадкой. 

    \item Внутреннюю поверхность пробки расширительного бачка радиатора после очистки от~коррозии смазать консервационной смазкой СХК или~ПВК. Закрыть жалюзи радиатора. 

    \item \begin{enumerate}
        \item Установить транспортер-тягач против лежней на~расстоянии 1~м от~них. Лежни должны быть на~1~м длиннее опорной поверхности гусениц. 
        \item Очистить гусеницы от~грязи, разъединить их и~проверить состояние траков и~пальцев, негодные~--- заменить. 
        \item Подготовить к~окраске и~окрасить лаком №177 свободные ветви гусениц и~необрезиненные ободы опорных катков, не~допуская при~этом окраски их резиновых бандажей и~попадание лака в~проушины звеньев гусениц. 
    \end{enumerate}
    \note{При~отсутствии лака №177 допускается применение смеси из~60\% битума и~40\% бензина.}
            

    \item Очистить от~коррозии все детали ходовой части (направляющие и~ведущие колеса, опорные катки, балансиры, кривошипы и~др.), частично или~полностью загрунтовать и~окрасить их нитроэмалью.

    \item \begin{enumerate}
        \item Очистить и~вымыть резиновые бандажи опорных катков и~соединить гусеницы. 
            
        \item Пустить и~прогреть двигатель на~режиме \range{1300}{1500}~об/мин. 
            
        \item Удалить воду из~бачка устройства для~смыва грязи с~ветровых стекол кабины путем полной ее выработки. Остаток воды с~осадком слить через сливную пробку, продуть бачок и~трубопроводы сжатым воздухом пневмосистемы транспортера-тягача. 
            
        \item Установить транспортер-тягач на~лежни, поставить рычаг переключения передач в~нейтральное положение и~остановить двигатель, перекрыв подачу топлива. 
            
        \item Отрегулировать натяжение гусениц, подготовить и~окрасить неокрашенные ее участки. 
    \end{enumerate}

    \item Слить отстой топлива в~количестве \range{5}{6}~л из~каждой группы топливных баков и~дозаправить их до~нормы. Заменять дизельное топливо толь ко~в~том случае, если оно не~соответствует времени года 

    \item \begin{enumerate}
        \item Произвести консервацию зеркальных поверхностей цилиндров двигателя, для~чего: снять форсунки, провернуть стартером коленчатый вал двигателя (\range{два}{три} включения стартера на \range{5}{6}~сек, с~интервалом между включениями \range{10}{15}~сек), установить поршень 1-го цилиндра в~положение НМТ и~последовательно, начиная с~1-го цилиндра, залить через отверстие стакана форсунки по \range{90}{100}~см$^3$ разогретого до \range{70}{80}\celc{} рабоче-консервационного масла во~все цилиндры. 
        
        \note{Одновременно масло можно заливать: в~1 и~6; 3 и~5; 4 и~7; 2 и~8 цилиндры.}
        
        \item Провернуть стартером коленчатый вал двигателя (\range{два}{три} включения по \range{5}{6}~сек, с~интервалом \range{10}{15}~сек), установить на~место форсунки. 
    \end{enumerate}
    \note{Рабоче-консервационное масло приготовляется непосредственно в~воинских частях путем добавления к~моторному или~трансмиссионному маслу 10\% присадки АКОР-1, подогретой до \range{60}{70}\celc{}, с~последующим тщательным перемешиванием.}

    \item Выполнить консервацию топливного насоса высокого давления и~регулятора числа оборотов путем заливки нагретого до~\range{70}{80}\celc{} рабоче-консервационного масла через отверстие для~указателя уровня масла (заливное отверстие) картера топливного насоса и~заливное отверстие картера регулятора оборотов при~рычаге ручной подачи топлива в~положении выключенной подачи. 
        
        После заполнения маслом картеров полностью, переместить несколько раз педаль управления или~рычаг ручной подачи топлива. Затем слить масло до~уровня меток на~контрольных щупах: в~картере регулятора~--- через сливное отверстие для~масла; в~топливном насосе~--- путем отсоса через отверстие под~маслоизмерительный щуп. 

    \item Выполнить работы по~консервации зеркальных поверхностей цилиндров компрессора: 
        \begin{enumerate}
            \item вывернуть пробки клапанов, вынуть пружины и~клапаны; 
            \item залить 20~см$^3$ рабоче-консервационного масла через отверстия клапанов;
            \item прокрутить шкив компрессора вручную;
            \item поставить на~место клапаны, пружины и~пробки клапанов. 
        \end{enumerate}
    \item Загерметизировать тканью ТТ и~замазкой ЗЗК-3У крышку воздухоотвода, трубку забора воздуха воздухофильтра, отверстие под~маслоизмерительный щуп, шахту выпуска, жалюзи радиатора, маслозаливную горловину двигателя и~заливную горловину масляного бака главной передачи. 

    \item Заполнить с~помощью топливоподкачивающего насоса топливоподводящие магистрали топливом, удалив из~них полностью воздух. 

    \item \begin{enumerate}
        \item Осмотреть электропровода, при~необходимости удалить с~их изоляции и~оплетки нефтепродукты (промыть бензином Б-70 и~протереть насухо). 
        
        \item Очистить наружные поверхности генератора и~стартера, при~необходимости подкрасить их. 
            
        \item Проверить затяжку всех зажимов электропроводки и~покрыть их поверхности тонким слоем лака №177. 
    \end{enumerate}
    

    \item \begin{enumerate}
        \item Проверить состояние осветительных и~светосигнальных приборов, сняв наружные ободки и~рассеиватели; очистить поверхность от~пыли, грязи и~продуктов коррозии; промыть рассеиватели с~обеих сторон и~протереть насухо. 
        
        \item При~обнаружении на~отражательной поверхности оптических элементов фар коррозии, пыли или~влаги разобрать их и~протереть чистой замшей или~фланелью. 
            
        \item Заменить неисправные рассеиватели и~уплотнительные прокладки во~всех приборах. 
            
        \item При~необходимости подкрасить внутренние и~наружные поверхности корпусов приборов и~наружных ободков. 
            
        \item Крышки СМУ фар опустить. 
    \end{enumerate}
    \note{Оптические элементы фар без~надобности не~разбирать. При~заметных повреждениях зеркальной поверхности оптические элементы заменить новыми. При~сборке осветительных приборов винты крепления наружных ободков затягивать равномерно со~всех сторон.}

    \item Снять аккумуляторные батареи, очистить их от~грязи, протереть, проверить и~при~необходимости отправить в~аккумуляторную для~обслуживания и~подзарядки. 
    Неисправные наконечники и~клеммы стартерных проводов заменить исправными, после чего смазать их консервационной смазкой СХК. 
    Поставить аккумуляторные батареи на~место. 

    \item Снять ремни привода вентилятора, генератора и~компрессора; очистить рабочие поверхности приводных шкивов от~продуктов коррозии и~окрасить. 
    Надеть ремни и~установить нормальное натяжение приводных ремней. 

    \item Прочистить и~загерметизировать сапуны картеров агрегатов силовой передачи, топливного насоса высокого давления, регулятора числа оборотов и~редуктора привода вентилятора. 

    \item Восстановить поврежденную окраску на~наружных поверхностях агрегатов и~узлов транспортера-тягача, при~необходимости полностью его окрасить. 

    \item Смазать консервационной смазкой СХК или~ПВК шарнирные соединения механизма запирания амбразур, механизма запирания заливных горловин, крышки люка отсека силовой передачи задних дверей кузова, неокрашенные поверхности тяговосцепного прибора и~шаровые опоры амбразур. 

    \item Свернуть и~уложить на~сиденья коврики пола кабины, предварительно вымыв и~высушив их. Пол тщательно очистить от~продуктов коррозии и~при~необходимости подкрасить. 
    Рычаги и~педали управления очистить, окрасить и~поставить в~нейтральное положение. 

    \item \begin{enumerate}
        \item Проверить по~комплектовочной ведомости, очистить, смазать рабочие поверхности консервационной смазкой СХК или~ПВК и~при~необходимости окрасить инструмент механика-водителя, возимый комплект запасных частей, шанцевый инструмент и~уложить их на~место. 
        \item Проверить вес заряда углекислоты в~огнетушителе и~при~не~обходимости зарядить его. 
    \end{enumerate}

    \item \begin{enumerate}
        \item При~необходимости очистить брезент от~пыли и~грязи, выполнить работы по~обработке (перепропитке) их химическим составом ПХС-55. 
        \item Укрыть транспортер-тягач брезентом. Для~предохранения брезента от~потертостей и~разрывов на~острые углы транспортера-тягача установить деревянные или~войлочные подкладки. 
        \item Брезент должен быть закреплен так, чтобы он не~соприкасался с~землей, не~раздувался ветром; на~поверхности брезента нежелательно скопление снега и~влаги. Можно транспортер-тягач укрыть чехлом-полотнищем, а~тент и~брезент сдать на~склад. 
    \end{enumerate}

    \item Закрыть смотровые окна, амбразуры и~заливные горловины внутренними запорами, а~люки трансмиссионного отделения, водителя, задние корпуса, крыши корпуса и~двигателя закрыть внешними запорами и~опломбировать.

\end{enumerate}


\section{Подготовка к~длительному хранению}
\subsubsection{Содержание работ и~технические условия}
\begin{enumerate}
    \item \begin{enumerate}
        \item 
        Тщательно очистить от~грязи и~вымыть транспортер-тягач, слить воду из~корпуса, протереть насухо наружные поверхности агрегатов и~узлов, из~труднодоступных мест воду удалить сжатым воздухом. 
        \item 
        При~проведении уборочно-моечных работ следить, чтобы не~попадала вода в~впускную и~выпускную системы двигателя, котел и~редуктор подогревателя и~на~приборы электрооборудования. 
        \item 
        Выполнить техническое обслуживание №2 и~дополнительные работы, предусмотренные через одно ТО-2 при~сезонном обслуживании, 
    \end{enumerate}

    \item \begin{enumerate}
        \item 
        Приготовить рабоче-консервационное масло путем добавления к~моторному маслу 10\% присадки АКОР-1, подогретой до \range{60}{70}\celc{}, с~последующим тщательным перемешиванием. 
        \item 
        Пустить двигатель, прогреть и~слить масло из~картера двигателя. Залить приготовленное рабоче-консервационное масло в~картер двигателя до~нормы. 
    \end{enumerate}
    \note{Для~приготовления рабоче-консервационного масла использовать только штатные зимние сорта моторных масел.}

    \item \begin{enumerate}
        \item 
        Промыть, протереть и~осмотреть дюритовые шланги. Шланги, имеющие расслоение с~торцов, вздутия и~широкие (до~1~мм) трещины на~поверхности, заменить. 
        \item 
        Очистить и~окрасить стяжные хомуты шлангов. Резьбовую часть стяжных болтов смазать консервационной смазкой СХК или~ПВК. 
    \end{enumerate}

    \item
    Заправить системы охлаждения двигателя водой или~низкозамерзающей жидкостью (антифризом) с~трехкомпонентной присадкой (по~0,05\% весовых нитрита натрия, тринатрия фосфата и~двухромокислого натрия), подогретыми зимой до \celc{50}. 

    \item
    Очистить от~коррозии внутреннюю поверхность пробки расширительного бачка, после очистки смазать консервационной смазкой СХК или~ПВК. Закрыть жалюзи радиатора. 

    \item \begin{enumerate}
        \item 
        Снять генератор с~двигателя и~очистить наружные поверхности; снять защитные ленты; осмотреть состояние щеток и~коллектора, обдуть их сжатым воздухом. 
        \item 
        При~подгорании (окислении) коллектора зачистить его стеклянной шкуркой №00, после чего обдуть сжатым воздухом и~протереть чистой ветошью, смоченной бензином Б-70. 
        \item 
        Установить генератор на~двигатель, предварительно при~необходимости подкрасить наружные поверхности. 
    \end{enumerate}

    \item \begin{enumerate}
        \item 
        Приготовить рабоче-консервационное масло путем добавления к~маслу МТ-14п или~МТ-16п присадки АКОР-1, подогретой до \range{60}{70}\celc{}, с~последующим тщательным перемешиванием. 
        \item 
        Заправить все агрегаты силовой передачи рабоче-консервационным маслом до~нормы. 
    \end{enumerate}

    \item \begin{enumerate}
        \item 
        Установить транспортер-тягач против лежней на~расстоянии 1~м от~них. Лежни должны быть на~1~м длиннее опорной поверхности гусениц. 
        \item 
        Очистить гусеницы от~грязи, разъединить их и~проверить состояние траков и~пальцев, негодные~--- заменить. 
        \item 
        Подготовить к~окраске и~окрасить лаком №177 свободные ветви гусениц и~необрезиненные ободы опорных катков, не~допуская при~этом окраски их резиновых бандажей и~попадание лака в~проушины звеньев гусениц. 

       \note{При~отсутствии лака №177 допускается применение смеси из~60\% битума и~40\% бензина.}
    \end{enumerate}

    \item \begin{enumerate}
        \item 
        Очистить от~коррозии все детали ходовой части (направляющие и~ведущие колеса, опорные катки, балансиры, кривошипы и~др.). 
        \item 
        Снять выборочно по~одному опорному катку с~каждой стороны и~осмотреть ступицы, подшипники и~уплотнения. При~обнаружении следов коррозии на~деталях осматриваемых катков снять все катки, проверить их состояние и~при~необходимости очистить детали катков от~продуктов коррозии. Если дефекты обнаружены не~будут, то можно ограничиться удалением старой смазки при~снятых колпаках катков, промывкой дизельным топливом и~заполнением ступиц опорных катков рабоче-консервационным маслом. 
        \item 
        Удалить старую смазку из~картеров бортовых передач, направляющих колес и~залить рабоче-консервационное масло до~уровня контрольных отверстий. 
        \item 
        Частично или~полностью окрасить детали ходовой части нитроэмалью, а~на~резиновые бандажи опорных катков нанести защитное покрытие СПО-46 (ПС-40 или~алюминиевую краску). 
    \end{enumerate}

    \item \begin{enumerate}
        \item 
        Соединить гусеницы, пустить и~прогреть двигатель на~режиме \range{1300}{1500}~об/мин. 
        \item 
        Удалить воду из~бачка устройства для~смыва грязи с~ветровых стекол кабины путем полной ее выработки. Остаток воды с~осадком слить через сливную пробку, продуть бачок и~трубопроводы сжатым воздухом пневмосистемы транспортера-тягача. 
        \item 
        Установить транспортер-тягач на~лежни, поставить рычаг переключения передач в~нейтральное положение и~остановить двигатель, перекрыв подачу топлива. Выключить выключатель массы. 
        \item 
        Отрегулировать натяжение гусениц, подготовить и~окрасить неокрашенные ее участки. 
    \end{enumerate}
    \note{Если консервация проводится без~передвижения (транспортер-тягач установлен на~лежни до~консервации), то для~создания масляной защитной пленки на~рабочих поверхностях деталей, не~соединяя гусениц, пустить двигатель, поработать в~течение \range{3}{5}~мин на~режиме \range{1300}{1500}~об/мин, затем прокрутить агрегаты силовой передачи и~выполнить оставшиеся операции, предусмотренные в~п.~11.}

    \item \begin{enumerate}
        \item 
        Слить топливо из~топливных баков, при~загрязнении топливные баки тщательно промыть; извлечь и~промыть фильтры грубой и~тонкой очистки, а~также заливных горловин топливных баков, промыть сапунные отверстия в~пробках и~продуть сжатым воздухом, установить фильтры на~место. 
        \item 
        Осмотреть топливопроводы и~проверить качество крепежных деталей, продуть топливопроводы сжатым воздухом; при~необходимости частично или~полностью окрасить все приборы топливной системы; топливные баки заправить зимним дизельным топливом (в~северных районах --- арктическим дизельным топливом), заполнить топливную систему топливом с~помощью топливоподкачивающего насоса, удалив из~нее полностью воздух. 
    \end{enumerate}

    \item \begin{enumerate}
        \item 
        Слить охлаждающую жидкость с~раствором трехкомпонентной присадки из~системы охлаждения двигателя. Слитую низкозамерзающую жидкость (антифриз) хранить в~соответствии с~действующим положением о~хранении ядовитых технических жидкостей. 
        \item 
        Сливные краники вывернуть, очистить от~осадков и~коррозии, внутренние и~наружные поверхности смазать консервационной смазкой СХК или~ПВК и~поставить на~место. Закрыть жалюзи радиатора. 
    \end{enumerate}

    \item \begin{enumerate}
        \item 
        Произвести консервацию зеркальных поверхностей цилиндров двигателя, для~чего: снять форсунки, провернуть стартером коленчатый вал двигателя (\range{два}{три} включения стартера по. \range{5}{6}~сек, с~интервалом между включениями \range{10}{15}~сек), установить поршень 1-го цилиндра в~положение НМТ и~последовательно, начиная с~1-го цилиндра, залить через отверстие стакана форсунки по \range{90}{100}~см$^3$ разогретого до \range{70}{80}\celc{} рабоче-консервационного масла во~все цилиндры. 
        
        \note{Одновременно масло можно заливать в~1 и~6; 3 и~5; 4 и~7; 2 и~8 цилиндры.}

        \item 
        Провернуть стартером коленчатый вал двигателя (\range{два}{три} включения по \range{5}{6}~сек, с~интервалом \range{10}{15}~сек) и~установить на~место форсунки. 
    \end{enumerate}

    \item \begin{enumerate}
        \item 
        Выполнить консервацию топливного насоса высокого давления и~регулятора числа оборотов путем заливки нагретого до \range{70}{80}\celc{} рабоче-консервационного масла через отверстие для~указателя уровня масла (заливное отверстие) картера топливного насоса и~заливное отверстие картера регулятора оборотов при~выключенном рычаге ручной подачи топлива. 
        \item 
        После заполнения маслом картеров полностью, переместить несколько раз педаль управления или~рычаг ручной подачи топлива. Затем слить масло из~картеров до~уровня меток на~контрольных щупах: в~картере регулятора~--- через сливное отверстие для~масла; в~топливном насосе~--- путем отсоса через отверстие под~маслоизмерительный щуп. 
    \end{enumerate}

    \item
    Выполнить работы по~консервации зеркальных поверхностей цилиндров компрессора: вывернуть пробки клапанов, вынуть пружины и~клапаны, залить по~20~см$^3$ рабоче-консервационного масла через отверстия клапанов, прокрутить шкив компрессора вручную и~поставить на~место, клапаны, пружины и~пробки клапанов. 

    \item
    Загерметизировать тканью ТТ и~замазкой ЗЗК-ЗУ крышку воздухоотвода, трубку забора воздуха воздухофильтра, отверстие под~маслоизмерительный щуп, шахту выпуска, жалюзи радиатора, маслозаливную горловину двигателя и~заливную горловину масляного бака главной передачи.

    \item
    Заполнить с~помощью топливоподкачивающего насоса топливоподводящие магистрали топливом, удалив из~них полностью воздух. 

    \item \begin{enumerate}
        \item 
        Проверить подогреватель системы охлаждения двигателя, вывернуть форсунку и~свечу накаливания и~смазать сопло форсунки и~спираль свечи консервационной смазкой СХК или~ПВК. 
        \item 
        Очистить при~необходимости котел подогревателя от~сажи и~продуктов коррозии. Окрасить снаружи частично или~полностью. 
    \end{enumerate}

    \item \begin{enumerate}
        \item 
        Осмотреть электропровода, при~необходимости удалить с~их изоляции и~оплетки нефтепродукты (промыть бензином Б-70 и~протереть насухо). 
        \item 
        Очистить наружные поверхности генератора и~стартера, при~необходимости подкрасить их. 
        \item 
        Проверить затяжку всех зажимов электропроводки и~покрыть их поверхности тонким слоем лака №177. 
    \end{enumerate}

    \item \begin{enumerate}
        \item 
        Проверить состояние осветительных и~светосигнальных приборов, сняв наружные ободки и~рассеиватели; очистить поверхности от~пыли, грязи и~продуктов коррозии; промыть рассеиватели с~обеих сторон и~протереть насухо. 
        \item 
        При~обнаружении на~отражательной поверхности оптических элементов фар коррозии, пыли или~влаги разобрать их и~протереть чистой замшей или~фланелью. 
        \item 
        Заменить неисправные рассеиватели и~уплотнительные прокладки во~всех приборах исправными. При~необходимости подкрасить внутренние и~наружные поверхности корпусов приборов и~наружных ободков. Крышки СМУ фар опустить. 
    \end{enumerate}
    \note{Оптические элементы фар без~надобности не~разбирать. При~заметных повреждениях зеркальной поверхности оптические элементы заменить новыми. При~сборке осветительных приборов винты крепления наружных ободков затягивать равномерно со~всех сторон.}

    \item \begin{enumerate}
        \item 
        Снять аккумуляторные батареи, очистить их от~грязи, протереть, проверить и~при~необходимости отправить в~аккумуляторную для~обслуживания и~подзарядки. 
        \item 
        Неисправные наконечники и~клеммы стартерных проводов заменить исправными, после чего смазать их консервационной смазкой СХК. Поставить аккумуляторные батареи на~место. 
    \end{enumerate}

    \item
    Снять ремни привода вентилятора, генератора и~компрессора; очистить рабочие поверхности приводных шкивов от~продуктов коррозии и~окрасить. Надеть ремни и~установить нормальное натяжение приводных ремней. 

    \item
    Прочистить и~загерметизировать сапуны картеров агрегатов силовой передачи, топливного насоса высокого давления, регулятора числа оборотов и~редуктора привода вентилятора. 

    \item
    Восстановить поврежденную окраску на~наружных поверхностях агрегатов и~узлов транспортера-тягача, при~необходимости полностью его окрасить. 

    \item
    Смазать консервационной смазкой СХК или~ПВК шарнирные соединения механизма запирания амбразур, механизма запирания заливных горловин, крышки люка отсека силовой передачи, задних дверей кузова, неокрашенные поверхности тяговосцепного прибора и~шаровые опоры амбразур. 

    \item \begin{enumerate}
        \item 
        Свернуть и~уложить на~сиденья коврики пола кабины, предварительно вымыв и~высушив их. 
        \item 
        Пол тщательно очистить от~продуктов коррозии и~при~необходимости подкрасить. 
        \item 
        Рычаги и~педали управления очистить, окрасить и~поставить в~нейтральное положение. 
    \end{enumerate}

    \item \begin{enumerate}
        \item 
        Проверить по~комплектовочной ведомости, очистить и~смазать рабочие поверхности инструмента консервационной смазкой СХК или~ПВК и~при~необходимости окрасить инструмент механика-водителя, возимый комплект запасных частей, шанцевый инструмент и~уложить их на~место. 
        \item 
        Проверить вес заряда углекислоты в~огнетушителе и~при~необходимости зарядить его. 
    \end{enumerate}

    \item \begin{enumerate}
        \item 
        Укрыть транспортер-тягач брезентом. Для~предохранения брезента от~потертостей и~разрывов на~острые углы транспортера-тягача установить деревянные или~войлочные подкладки. 
        \item 
        Брезент должен быть закреплен так, чтобы он не~соприкасался с~землей, не~раздувался ветром; на~поверхности брезента нежелательно скопление снега и~влаги. Можно транспортер-тягач укрыть чехлом-полотнищем, а~тент и~брезент сдать на~склад. 
    \end{enumerate}

    \item \begin{enumerate}
        \item 
        Закрыть смотровые окна, амбразуры и~заливные горловины внутренними запорами, а~люки отсека силовой передачи, водителя, задние корпуса, крышки корпуса и~двигателя закрыть внешними запорами и~опломбировать. 
        \item 
        Закрыть задние двери кузова, левую дверь кабины внешними запорами и~оп\-ломбировать. 
        \item 
        Опломбировать также верхний клапан тента и~заливные горловины. 
\end{enumerate}

\end{enumerate}



\section{Обслуживание транспортера-тягача, находящегося в~консервации} 
\subsection{Один раз в~месяц}
\subsubsection{Содержание работ и~технические условия}
\begin{enumerate}
    \item Проверить целостность всех пломб. 
    \item Убрать место стоянки транспортера-тягача и~очистить защитный (укрывочный) брезент от~пыли и~осадков. 
    \item Проверить положение транспортера-тягача на~лежнях и~устранить обнаруженные недостатки. 
\end{enumerate}
\note{В~зависимости от~погоды (обильные дожди, снегопады, ветры с~пылью) и~условий хранения транспортера-тягача может быть назначена внеочередная уборка с~определением необходимого объема работ.}



\subsection{Один раз в~шесть месяцев}
\subsubsection{Содержание работ и~технические условия}
\begin{enumerate}
    \item Выполнить на~транспортере-тягаче работы, проводимые один раз в~месяц. 
    \item Снять защитный брезент. 
    \item Распломбировать и~открыть в~сухую без~осадков погоду люки, для~проветривания; просушить брезент и~коврики. 
    \item Проверить состояние наружных поверхностей агрегатов и~механизмов; места, подвергшиеся коррозии, очистить, окрасить и~смазать консервационной смазкой СХК или~ПВК. 
    \item Проверить, не~подтекает~ли топливо, масло, охлаждающая и~амортизаторная жидкости; обнаруженные неисправности устранить. 
    \item Проверить и~при~необходимости восстановить герметизирующие оклейки на~узлах и~агрегатах. 
    \item Заменить в~агрегатах и~механизмах транспортеров-тягачей, находящихся на~кратковременной консервации, масло и~топливо сортами, соответствующими предстоящему периоду эксплуатации. 
    \item Проверить вес заряда в~углекислотном огнетушителе и~при~необходимости зарядить. 
    \item Закрыть и~опломбировать вскрытые при~обслуживании люки. 
    \item Укрыть транспортер-тягач защитным брезентом. 
\end{enumerate}


\subsection{Один раз в~год}
\subsubsection{Содержание работ и~технические условия}
\begin{enumerate}
    \item Выполнить на~транспортере-тягаче работы, проводимые один раз в~шесть месяцев. 
    \item Слить отстой дизельного топлива в~количестве \range{5}{6}~л из~каждой группы топливных баков. 
\end{enumerate}

Кроме этого, в~процессе хранения при~очередных годовых обслуживаниях на~всех транспортерах-тягачах длительной консервации один раз в~пять лет выполнить следующие работы: 
\begin{enumerate}
    \item Заменить масло в~картерах двигателя, топливного насоса высокого давления, регулятора числа оборотов и~редуктора привода вентилятора. 
    \item Заменить масло в~агрегатах силовой передачи. 
    \item Заменить дизельное топливо в~топливных баках. 
    \item Заменить амортизаторную жидкость в~амортизаторах. 
    \item Заменить масло в~опорных катках и~направляющих колесах. 
    \item Проверить состояние и~при~необходимости заменить дюритовые шланги, электропроводку и~резино-технические изделия. 
\end{enumerate}



\section{Снятие с~консервации}
\subsubsection{Содержание работ и~технические условия}
Работы первой очереди, выполняемые при~ускоренном снятии с~хранения: 
\begin{enumerate}
    \item Снять защитный брезент. 
    \item Распломбировать люки. 
    \item Установить аккумуляторные батареи (если они были сняты) в~гнезда, протереть клеммы и~зажимы проводов ветошью, смоченной бензином Б-70, закрепить батареи, присоединить к~ним провода, смазать клеммы пресс-солидолом и~проверить наличие тока в~электрических цепях. 
    \item \begin{enumerate} 
            \item Подготовить двигатель к~пуску. Проверить уровень масла в~картере и~наличие топлива в~баках. На~машинах длительной консервации, а~в~зимнее время и~на~машинах кратковременной консервации~--- заправить систему охлаждения водой или~низкозамерзающей жидкостью (зимой~--- низкозамерзающей жидкостью или~горячей водой, с~использованием подогревателя двигателя). 
            \item Разгерметизировать воздухоочиститель двигателя, выпускную шахту и~жалюзи радиатора. 
        \end{enumerate}
    \item Провернуть коленчатый вал двигателя на \range{10}{15} оборотов. 
    \item Пустить двигатель, прогреть его и~проверить работу на~всех режимах. 
    \item Проверить работу контрольно-измерительных приборов. 
    \item Вывести транспортер-тягач с~лежней. 
\end{enumerate}

Работы второй очереди, выполняемые после снятия с~хранения в~районах сосредоточения, на~привалах или~остановках: 
\begin{enumerate}
    \item Свернуть, уложить на~место и~закрепить защитный брезент. 
    \item Разгерметизировать сапуны картеров агрегатов силовой передачи, топливного насоса высокого давления, регулятора числа оборотов и~редуктора привода вентилятора. 
    \item Уложить коврики на~пол кузова МТЛБ транспортера-тягача. 
    \item Очистить инструмент от~слоя консервационной смазки и~уложить его на~место.
\end{enumerate}


